\PassOptionsToPackage{fleqn}{amsmath}
\documentclass[final,t]{beamer}
\usepackage[orientation=portrait,size=a1,scale=1.12]{beamerposter}

\usetheme{default}
\usecolortheme{seagull}
\usefonttheme{professionalfonts}

\usepackage[T1]{fontenc}
\usepackage{graphicx}
\usepackage{url}
\usepackage{newtxtext,newtxmath}
\usepackage{xcolor}
\usepackage{amsmath,amssymb}

\usepackage{hyperref}
\hypersetup{colorlinks=true,linkcolor=blue,urlcolor=blue,citecolor=blue}
\renewcommand\UrlFont{\color{blue}\rmfamily}
\urlstyle{rm}

\title{\parbox{0.93\linewidth}{\centering
Inverting the Formalization Workflow:\\[0.2em]
\Large Prototyping an MPC Protocol in Rocq with an LLM Agent}}
\author{Cheng-Hui Weng \quad Nagoya University}
\date{}

\newcommand{\kw}[1]{\textbf{#1}}
\newcommand{\todo}[1]{\textcolor{red}{#1}}

\begin{document}
\begin{frame}[t]
\vspace{-0.5em}
\begin{columns}[t,totalwidth=\linewidth]
  \begin{column}{0.49\linewidth}
    \begin{block}{Motivation}
      Formalization usually comes last:
      \[
      \text{learn} \to \text{design} \to \text{formalize}.
      \]
      This is costly when a protocol depends on several unfamiliar domains.

      \vspace{0.4em}
      We reverse the order and treat formalization as a prototyping medium:
      \[
      \text{idea} \to \text{formalize} \to \text{learn}.
      \]
    \end{block}

    \begin{block}{Inverted Workflow}
      \begin{itemize}
        \item Human proposes hypotheses, steers abstractions, and sets axiomatization boundaries.
        \item LLM drafts definitions, lemma statements, and candidate proof structure.
        \item Rocq/MathComp checks logical soundness at each step.
        \item Failed obligations reveal missing assumptions early.
      \end{itemize}
      \vspace{0.3em}
      \kw{Key loop:} draft $\to$ check $\to$ diagnose $\to$ refine.
    \end{block}

    \begin{block}{Case Study: SMC-PGG}
      Secure multi-party computation via covering-space structure:
      \begin{itemize}
        \item Monodromy representation \(\rho : G \to S_N\) acts on \(N\) sheets.
        \item Secret is encoded in fiber structure over a basepoint.
        \item Session-typed protocol layer implemented in Rocq.
      \end{itemize}
      The development is organized as protocol, search-space, reconstruction, and security layers.
    \end{block}
  \end{column}

  \begin{column}{0.49\linewidth}
    \begin{block}{Technical Highlights}
      \begin{itemize}
        \item \kw{Search-space analysis:} RAAG trace equivalence for commuting-action reordering.
        \item \kw{Combinatorics:} Cartier-Foata theorem used for exact trace counting via clique-polynomial recurrence.
        \item \kw{Reconstruction:} threshold interface instantiated through linear-code secret sharing.
        \item \kw{Security bound:}
          \[
          \mathrm{TV} \le \varepsilon + \frac{2(T-1)}{N}.
          \]
        \item \kw{Structure result:} abelian collapse in regular abelian settings; non-abelian structure is required for non-trivial security.
      \end{itemize}
    \end{block}

    \begin{block}{Current Scope and Status}
      \begin{itemize}
        \item Ongoing formalization: \kw{42 Rocq files}.
        \item Domains touched: group theory, combinatorics, algebraic geometry, coding theory, information theory.
        \item Remaining axioms are concentrated at the algebraic-geometry interface.
      \end{itemize}
    \end{block}

    \begin{block}{Methodological Takeaways}
      \begin{itemize}
        \item LLMs are useful not only for proof search, but also for cross-domain exploration.
        \item The prover supplies certainty; the human supplies judgment.
        \item Formalization can be an engine for learning, not only final verification.
      \end{itemize}
    \end{block}

    \begin{block}{QR / Links / Contact}
      \small
      Paper draft and artifacts:\quad \url{https://github.com/<org>/<repo>}

      \vspace{0.25em}
      Contact:\quad \href{mailto:cheng-hui.weng@example.com}{cheng-hui.weng@example.com}

      \vspace{0.4em}
      \todo{Replace URL, email, and add QR image if needed.}
    \end{block}
  \end{column}
\end{columns}
\end{frame}
\end{document}
